% Generated draft: requires mathematical and source review.
\documentclass[11pt]{article}
\usepackage[margin=25mm]{geometry}
\usepackage{fontspec}
\setmainfont{Latin Modern Roman}
\usepackage{amsmath,amssymb,mathtools,unicode-math}
\setmathfont{Latin Modern Math}
\usepackage{enumitem,needspace,xurl,hyperref}
\setlength{\parindent}{0pt}
\setlength{\parskip}{5pt}
\setlength{\emergencystretch}{3em}
\newcommand{\N}{\mathbb{N}}
\newcommand{\Z}{\mathbb{Z}}
\newcommand{\Q}{\mathbb{Q}}
\newcommand{\R}{\mathbb{R}}
\newcommand{\C}{\mathbb{C}}
\newcommand{\F}{\mathbb{F}}
\newcommand{\A}{\mathbb{A}}
\newcommand{\PP}{\mathbb{P}}
\newcommand{\E}{\mathbb{E}}
\newcommand{\eps}{\varepsilon}
\newcommand{\dd}{\,\mathrm{d}}
\begin{document}
\section*{UnsolvedMath ID: 800025}
\textbf{Source problem code:} catalogue entry
\subsection{Definitions and mathematical statement}
For an algebra \(A\) in a fixed signature, let \(\operatorname{End}(A)\) be the semigroup of all homomorphisms \(A\to A\), with multiplication \(fg=f\circ g\). Let \(\operatorname{Aut}(A)\) be the group of invertible endomorphisms. For a semigroup, these homomorphisms preserve multiplication.

Let \(\mathcal K\) be a class of algebras of a fixed signature. An algebra \(F_{\mathcal K}(X)\in\mathcal K\), with a map \(\iota:X\to F_{\mathcal K}(X)\), is free on \(X\) if every map \(u:X\to B\), for \(B\in\mathcal K\), extends uniquely to a homomorphism \(\bar u:F_{\mathcal K}(X)\to B\) satisfying \(\bar u\circ\iota=u\). This notation is used only where such a free algebra exists.

The explicitly named semigroup classes are the following. An idempotent semigroup, also called a band, satisfies \(s^2=s\) for every element \(s\). An inverse semigroup is a semigroup in which every \(s\) has a unique inverse \(t\) satisfying \(sts=s\) and \(tst=t\). A Clifford semigroup is an inverse semigroup whose idempotents are central: \(es=se\) for every idempotent \(e\) and every element \(s\). A commutative Clifford semigroup additionally satisfies \(st=ts\) for all elements \(s,t\).

\textbf{Problem (G. Mashevitzky; supplied source, Problem 26).} Describe the endomorphism semigroup and the tower of endomorphisms (automorphisms) for endomorphism semigroups of free semigroups in basic classes of semigroups, particularly the classes of idempotent semigroups, commutative Clifford semigroups, and Clifford semigroups. Seek the analogous descriptions for endomorphism semigroups of algebras free in basic classes of mathematical structures.

More explicitly, the literal nested reading starts with
\[
E=\operatorname{End}\bigl(F_{\mathcal K}(X)\bigr)
\]
and asks for a description of \(\operatorname{End}(E)\) and the endomorphism and automorphism towers based at \(E\). Possible finite-stage conventions are
\[
E_0=E,\qquad E_{n+1}=\operatorname{End}(E_n),
\]
and
\[
G_0=E,\qquad G_{n+1}=\operatorname{Aut}(G_n),
\qquad n\in\{0,1,2,\ldots\}.
\]
Here each \(E_n\) is treated as a semigroup; \(G_1\) is the automorphism group of the semigroup \(E\), and later automorphism stages are groups.

\textbf{Unresolved qualifications.} These recursions explain a possible meaning of the requested towers; the supplied record does not define them. In particular, it does not settle whether the intended starting object is \(F_{\mathcal K}(X)\) or its endomorphism semigroup \(E\), whether connecting homomorphisms are part of the requested description, or whether transfinite stages are intended. No limit-stage construction is specified. The rank \(|X|\), treatment of the empty generating set, remaining ``basic classes,'' signatures of the general algebras, and required form of a ``description'' are also unspecified. Thus the source gives a research programme rather than a uniquely quantified classification statement; no additional restrictions or universal rank claim are imposed here.

The supplied source reports that descriptions of endomorphism semigroups for the classes of all semigroups and of inverse semigroups were obtained in references labelled [MS] and [MSZ], and that automorphism towers of free inverse semigroups were described in [Mash3]. These are supplied background claims, not independently verified results in this draft. The local registry labels this problem open.
\subsection{Short English statement}
Describe endomorphism semigroups and iterated endomorphism or automorphism constructions associated with free semigroups, especially bands and Clifford semigroups, and extend the investigation to free algebras in other classes. The source leaves the ranks, broader classes, and precise tower conventions unspecified, including whether iteration begins with a free algebra or its endomorphism semigroup.
\subsection{Sources}
\begin{itemize}
\item Supplied local UnsolvedMath registry record, permanent ID \(800025\), exact source code \texttt{catalogue entry}, title ``The Lyapin's Notebook --- Problem 26.'' The registry records status \texttt{open}, lists Bershadsky, Kublanovsky, Mashevitzky in its proposer field, and gives proposed year 2026. These are supplied registry labels and attribution.
\item Bershadsky, Kublanovsky, Mashevitzky, \emph{The Lyapin's Notebook: Unsolved Problems in Semigroup Theory} (2026), Problem 26, attributed within the supplied statement to G. Mashevitzky. Supplied source URL: \url{https://arxiv.org/abs/2604.04763}. The provenance records access on 2026-10-16 and extraction by \texttt{arxiv\_latex}; these are supplied provenance details, not a claim of independent source access.
\item References [MS] and [MSZ], cited in the supplied statement for descriptions of endomorphism semigroups in the classes of all semigroups and inverse semigroups. Full bibliographic entries and precise hypotheses were not supplied.
\item Reference [Mash3], cited in the supplied statement for automorphism towers of free inverse semigroups. Its full bibliographic entry and precise hypotheses were not supplied.
\end{itemize}
\end{document}
